\section{Participants, informants, and evaluators}
\label{sec:role-distinction}

The central distinction is between producing data about oneself and applying expertise to a target object. A participant contributes evidence about a population of people, a psychological process, a distribution of acceptability, a demographic contrast, or an intervention effect. An evaluator contributes an assessment of an object: an essay, a translation, a diagnosis, a code, a performance, or a linguistic example.

This distinction is functional rather than honorific. Calling someone an evaluator does not make the judgment infallible, unregulated, or methodologically private. It changes the evidential question. The relevant issues become qualification, calibration, independence, adjudication, documentation, and conflict of interest, not sampling, participant burden, demographic representativeness, or the protection of private personal data.

For grammaticality work, the distinction matters because many expert judgments do not estimate speaker-preference distributions. The task is to decide whether a proposed sentence, construction, contrast, or diagnostic behaves as the theory says it does. That use of expertise is evaluative.
